\section{A closure engine for mathematics}
\label{sec:packaging-engine}

This section states the procedure that the exhibits apply.
It does not re-found analysis; it makes explicit the commitments that are usually left tacit when one passes from a discrete protocol to a packaged law.
There are five of them---staging, a ledger, a packaging contract, a descent contract, and route diagnostics---and Figure~\ref{fig:engine} shows how they fit together.

\begin{figure}[htbp]
\centering
\begin{tikzpicture}[
  font=\small,
  box/.style={draw=black!55, rounded corners=3pt, fill=black!3, align=center, inner sep=5pt, minimum height=11mm, text width=24mm},
  arr/.style={-{Latex[length=2.2mm]}, draw=black!65, line width=0.6pt},
  tag/.style={font=\scriptsize\bfseries, text=black!65}
]
\node[box] (stage) {staged family\\ $z_h$, \ $h\downarrow 0$};
\node[box, right=6mm of stage] (ledger) {defect ledger\\ $\Def(h)$};
\node[box, right=6mm of ledger] (pack) {package\\ class $[\{z_h\}]$};
\node[box, right=6mm of pack] (law) {packaged law\\ $F^\sharp$ on classes};
\node[box, below=9mm of ledger, text width=60mm] (routes) {two routes $R_1(h),R_2(h)$ to the same object\\ route mismatch $\RM(h)$};
\draw[arr] (stage) -- (ledger);
\draw[arr] (ledger) -- node[above, font=\scriptsize]{shrinks?} (pack);
\draw[arr] (pack) -- node[above, font=\scriptsize]{descends?} (law);
\draw[arr] (stage.south) |- (routes.west);
\draw[arr] (routes.east) -| (pack.south);
\node[tag, above=0.5mm of stage] {P4 staging};
\node[tag, above=0.5mm of ledger] {P6 accounting};
\node[tag, above=0.5mm of pack] {P5 packaging};
\node[tag, above=0.5mm of law] {P1, P2 descent};
\node[tag, below=0.5mm of routes] {P3 protocol};
\end{tikzpicture}
\caption{The closure engine. A staged family is audited by a defect ledger. If the defect vanishes under refinement in a genuine pseudometric, the family defines a packaged object (Proposition~\ref{prop:packaging}). A discrete operator becomes a packaged law only if it respects the packaging (Proposition~\ref{prop:descent}), which may require rescaling or rewriting it. Independently, two routes to the same object are compared by a route mismatch that should shrink in stable regimes.}
\label{fig:engine}
\end{figure}

\subsection{Staging (\Pfour)}
Fix a family of discrete protocol spaces $\{Z_h\}_{h>0}$ indexed by a refinement parameter: a mesh size, a truncation cutoff, a partition diameter.
A point $z_h\in Z_h$ is a description at stage $h$.
Often there are comparison maps $\rho_{h\to h'}:Z_h\to Z_{h'}$ for $h'<h$ that carry coarse objects to finer resolution (interpolation, refinement of a partition, extension of a truncation).
Staging is what turns ``infinitely many small steps'' into a question about the behaviour of a family as $h\downarrow 0$, and therefore into something that can be tested.

\subsection{Ledgers and feasibility gates (\Psix)}
A closure claim is meaningful only relative to a ledger of nonnegative diagnostics that measure what a compression loses.
We use two kinds:
\begin{enumerate}
  \item \emph{approximation errors}, comparing a candidate quantity across refinement, $\|Q_h-Q_{h'}\|$, or against a reference, $\|Q_h-Q_\star\|$;
  \item \emph{coherence defects}, measuring the failure of an intended algebraic constraint, such as a Leibniz defect for a product rule or a route mismatch between two protocols.
\end{enumerate}
A \emph{feasibility gate} is a requirement on the ledger, either asymptotic ($\Def(h)\to 0$ as $h\downarrow 0$) or, in finite experiments, a contraction test such as $\Def(h')\le\theta\,\Def(h)$ for $h'<h$ and some $\theta<1$, together with an absolute threshold.
A finite gate checks finitely many stages; it is evidence about the trend, not a proof of the limit.

\subsection{Packaging as identification (\Pfive)}
\label{subsec:packaging}
Packaging turns a staged family into an object by identifying families whose differences vanish under refinement.
This is the pattern of the Cauchy completion of the rationals and of identifying two difference quotients that differ by $o(1)$.
For it to define objects at all, the identification must be an equivalence relation, and an arbitrary nonnegative defect does not guarantee that.

\begin{example}[A defect that does not package]\label{ex:nontransitive}
On three points $a,b,c$ put $d(a,b)=d(b,c)=0$, $d(a,c)=1$, with $d$ symmetric and zero on the diagonal.
Taking $d$ independent of the stage, ``$d$ tends to zero'' relates $a$ to $b$ and $b$ to $c$ but not $a$ to $c$.
\end{example}

The missing ingredient is the triangle inequality.
We therefore make the comparison space part of the packaging contract.

\begin{proposition}[Packaging contract]\label{prop:packaging}
Let $(X,d)$ be a pseudometric space. For sequences $a,b:\mathbb N\to X$ write $a\sim b$ when $d(a_n,b_n)\to 0$. Then $\sim$ is an equivalence relation.
\end{proposition}

\begin{proof}
Reflexivity and symmetry follow from $d(x,x)=0$ and $d(x,y)=d(y,x)$. For transitivity, $0\le d(a_n,c_n)\le d(a_n,b_n)+d(b_n,c_n)\to 0$.
\end{proof}

\noindent In Lean this is \leanname{Closure.vanishingDistSetoid}.
Several ledger coordinates can be combined through a joint pseudometric or as a conjunction of such relations.
Three further points keep the contract honest.
\begin{itemize}
  \item \emph{A quotient is not a limit.} The quotient by $\sim$ of \emph{all} sequences is not a completion: bounded oscillating scalar sequences survive in $\ell^\infty/c_0$ without representing any real number. A completion requires restricting to Cauchy sequences, and convergence in an existing complete space must be proved there.
  \item \emph{Algebra need not descend.} On unbounded scalar sequences, multiplication does not respect $\sim$: $n$ and $n+1/n$ are equivalent, but their squares differ by $2+1/n^2$. On uniformly bounded families it does, by $|ab-a'b'|\le|a|\,|b-b'|+|b'|\,|a-a'|$.
  \item \emph{The stage dependence must be uniform.} A comparison that changes with the stage needs the same triangle-type control uniformly in the stage; naming a ledger does not supply it.
\end{itemize}

\subsection{Descent of operators (\Pone, \Ptwo)}
A discrete operator family $F_n$ becomes a packaged law only if equivalent inputs give equivalent outputs.
A uniform Lipschitz estimate is a sufficient condition.

\begin{proposition}[Descent contract]\label{prop:descent}
Let $X,Y$ be pseudometric spaces and $F_n:X\to Y$ maps with $d_Y(F_n x,F_n y)\le K\,d_X(x,y)$ for all $n,x,y$, where $K\ge 0$ does not depend on $n$.
If $d_X(a_n,b_n)\to 0$, then $d_Y(F_n a_n,F_n b_n)\to 0$. Hence $(a_n)\mapsto(F_n a_n)$ induces a map on equivalence classes.
\end{proposition}

\begin{proof}
$0\le d_Y(F_n a_n,F_n b_n)\le K\,d_X(a_n,b_n)\to 0$.
\end{proof}

\noindent In Lean this is \leanname{Closure.stage_operator_respects_vanishing_dist}.
The condition is sufficient, not necessary, and it genuinely fails for the operator at the heart of calculus if the comparison is chosen carelessly.

\begin{example}[Difference quotients do not descend on uniform values]\label{ex:no-descent}
Let $h_n=1/n$, $f_n(x)=h_n\sin(x/h_n)$ and $g_n=0$. Then $\|f_n-g_n\|_\infty\le h_n\to 0$, yet $\delta_{h_n}f_n(0)=\sin(1)\ne 0=\delta_{h_n}g_n(0)$.
The Lipschitz constant of $\delta_h$ with respect to the uniform norm is of order $1/h$, so Proposition~\ref{prop:descent} does not apply.
\end{example}

The repair is a richer comparison, which is an instance of \Pone.
If $f,g\in C^1$, the mean value theorem gives $\|\delta_h(f-g)\|_\infty\le\|f'-g'\|_\infty$ on any interior domain, so equivalence of inputs in the $C^1$ norm descends to equivalence of outputs in the $C^0$ norm.
For an operator that maps a function algebra to itself, one works on a derivative-closed class such as polynomials or smooth functions with derivative seminorms.

\subsection{Route mismatch (\Pthree)}
When two reasonable routes $R_1(h)$ and $R_2(h)$ produce outputs on a common space, the relative mismatch~\eqref{eq:rm-generic} is a computable defect, and its decay under refinement indicates a stable regime.
For bounded linear routes the absolute version of this defect has a clean meaning: it controls how much the direct route can amplify compared with the composite one.

\begin{proposition}[Route defect bounds gain]\label{prop:route-gain}
Let $A$, $B$, $C$ be continuous linear maps between normed spaces with $C$ and $A\circ B$ of the same type. Then
\[
  \|C\|\;\le\;\|A\|\,\|B\|+\|C-A\circ B\|.
\]
For a continuous linear endomap $L$, moreover $L\circ L-L=L\circ(L-\mathrm{Id})$, so $\|L\circ L-L\|\le\|L\|\,\|L-\mathrm{Id}\|$.
\end{proposition}

\noindent These are \leanname{Closure.route_gain_bound}, \leanname{Closure.idempotence_factorization}, and \leanname{Closure.idempotence_defect_bound}.
Linearity matters: for $E(x)=|x|$ on $\mathbb R$ one has $E\circ E-E=0$ but $E\bigl((E-\mathrm{Id})(-1)\bigr)=2$, so the factorization fails for nonlinear maps.
The bound concerns the absolute operator-norm defect.
The relative diagnostic used in the exhibits is evaluated on finitely many sampled inputs; rescaling it recovers only the absolute discrepancy on those inputs, not the operator-norm defect. (For $A=B=\mathrm{Id}$ and $C=\mathrm{diag}(1,M)$ with $M>1$ on Euclidean $\mathbb R^2$, sampling only $e_1$ shows zero mismatch although $\|C-A\circ B\|=M-1$.) Applying the gain bound therefore requires a separately established operator-norm estimate, and even a zero route defect at each step does not by itself control growth over many steps.
Above all, route mismatch is a diagnostic of closure feasibility; it does not certify a direction of time or any irreversibility.

\subsection{Closure ladders}
The engine composes.
Once a packaged layer is stable it becomes the substrate for further staging and packaging: the reals arise by completing the rationals, calculus arises by packaging difference and sum protocols under coherence constraints, and geometric obstructions appear as route dependence when local descriptions are patched together.
The exhibits that follow are small, auditable instances of this pattern.
